\begin{center}
    \textbf{\LARGE ĐỀ KIỂM TRA ĐÁNH GIÁ GIỮA KÌ I}\\
    \textbf{\Large MÔN: HÓA HỌC 10 -- BỘ SÁCH CÁNH DIỀU / KNTT / CHÂN TRỜI}\\
    \textbf{\large NỘI DUNG: NGUYÊN TỬ -- BẢNG TUẦN HOÀN CÁC NGUYÊN TỐ HÓA HỌC -- ĐỀ SỐ 01}\\
    \textit{Thời gian làm bài: 45 phút, không kể thời gian phát đề}
\end{center}
\hrule
\vspace{0.5cm}

\noindent\textbf{\large PHẦN I. Câu trắc nghiệm nhiều phương án lựa chọn.} Thí sinh trả lời từ câu 1 đến câu 12. Mỗi câu hỏi thí sinh chỉ chọn một phương án.

\begin{ex}
    Đối tượng nghiên cứu của hóa học là
    \choice
    {các hiện tượng tự nhiên như sấm sét, bão lũ, động đất, núi lửa phun trào}
    {\True thành phần, cấu trúc, tính chất và sự biến đổi của chất cùng các ứng dụng của chúng}
    {sự chuyển động cơ học của các thiên thể và vật thể trong không gian vũ trụ}
    {các quy luật sinh trưởng, sinh sản và phát triển của các loài sinh vật}
    \loigiai{
        Hóa học là ngành khoa học tự nhiên nghiên cứu về thành phần, cấu trúc, tính chất và sự biến đổi của chất cũng như ứng dụng của chúng trong đời sống và sản xuất.
    }
\end{ex}

\begin{ex}
    Ứng dụng nào sau đây thể hiện rõ nét vai trò của hóa học đối với sự phát triển kinh tế nông nghiệp?
    \choice
    {Sản xuất màn hình hiển thị điện tử và chip bán dẫn}
    {Xây dựng mạng lưới định vị vệ tinh và viễn thông toàn cầu}
    {\True Sản xuất phân bón hóa học và thuốc bảo vệ thực vật nhằm nâng cao năng suất cây trồng}
    {Quy hoạch hệ thống giao thông đường bộ và phân luồng xe cộ}
    \loigiai{
        Hóa học đóng vai trò quan trọng trong nông nghiệp thông qua việc nghiên cứu, sản xuất các loại phân bón hóa học (đạm, lân, kali, phân vi lượng), thuốc bảo vệ thực vật, chất kích thích sinh trưởng giúp cây trồng phát triển tốt và nâng cao năng suất mùa màng.
    }
\end{ex}

\begin{ex}
    Để nghiên cứu một vấn đề hóa học theo phương pháp tìm hiểu thế giới tự nhiên dưới góc độ hóa học, trình tự các bước thực hiện hợp lí là
    \choice
    {\True Đề xuất vấn đề $\rightarrow$ Đưa ra giả thuyết khoa học $\rightarrow$ Lập kế hoạch và thực hiện nghiên cứu $\rightarrow$ Thu thập, xử lí dữ liệu $\rightarrow$ Kết luận}
    {Lập kế hoạch và thực hiện nghiên cứu $\rightarrow$ Đề xuất vấn đề $\rightarrow$ Thu thập dữ liệu $\rightarrow$ Đưa ra giả thuyết $\rightarrow$ Kết luận}
    {Đưa ra giả thuyết khoa học $\rightarrow$ Kết luận $\rightarrow$ Đề xuất vấn đề $\rightarrow$ Thực hiện nghiên cứu $\rightarrow$ Thu thập dữ liệu}
    {Thu thập dữ liệu $\rightarrow$ Đề xuất vấn đề $\rightarrow$ Đưa ra giả thuyết khoa học $\rightarrow$ Thực hiện nghiên cứu $\rightarrow$ Kết luận}
    \loigiai{
        Phương pháp tìm hiểu thế giới tự nhiên dưới góc độ hóa học gồm các bước cơ bản:
        \begin{itemize}
            \item[(1)] Đề xuất vấn đề nghiên cứu (quan sát hiện tượng, đặt câu hỏi);
            \item[(2)] Đưa ra giả thuyết khoa học;
            \item[(3)] Lập kế hoạch và tiến hành thực nghiệm (thí nghiệm kiểm chứng, thu thập thông tin);
            \item[(4)] Thu thập, phân tích và xử lí dữ liệu;
            \item[(5)] Thảo luận và rút ra kết luận.
        \end{itemize}
    }
\end{ex}

\begin{ex}
    Thành phần cấu tạo của hầu hết các nguyên tử gồm
    \choice
    {hạt nhân chứa electron và lớp vỏ chứa proton}
    {hạt nhân chứa neutron và lớp vỏ chứa proton}
    {\True hạt nhân gồm proton, neutron và lớp vỏ gồm các electron}
    {hạt nhân gồm proton, electron và lớp vỏ gồm các neutron}
    \loigiai{
        Hầu hết các nguyên tử đều được tạo nên từ hai phần: hạt nhân ở tâm gồm các hạt proton và neutron (riêng nguyên tử hydrogen đồng vị $^1_1\text{H}$ chỉ có proton, không có neutron) và lớp vỏ nguyên tử gồm các electron mang điện tích âm chuyển động xung quanh hạt nhân.
    }
\end{ex}

\begin{ex}
    Trong nguyên tử, loại hạt cơ bản mang điện tích dương và có khối lượng xấp xỉ bằng $1\text{ amu}$ là
    \choice
    {electron}
    {\True proton}
    {neutron}
    {photon}
    \loigiai{
        - Hạt proton mang điện tích $+1$ ($+1{,}602 \times 10^{-19}\text{ C}$) và có khối lượng $m_p \approx 1\text{ amu}$.\\
        - Hạt neutron không mang điện và có khối lượng $m_n \approx 1\text{ amu}$.\\
        - Hạt electron mang điện tích $-1$ ($-1{,}602 \times 10^{-19}\text{ C}$) và có khối lượng rất nhỏ $m_e \approx 0{,}00055\text{ amu}$.
    }
\end{ex}

\begin{ex}
    Nguyên tố hóa học là tập hợp các nguyên tử có cùng
    \choice
    {số neutron trong hạt nhân}
    {số khối $A$}
    {\True số proton trong hạt nhân (điện tích hạt nhân)}
    {tổng số hạt proton và neutron}
    \loigiai{
        Nguyên tố hóa học là tập hợp những nguyên tử có cùng điện tích hạt nhân (tức là có cùng số proton trong hạt nhân hay cùng số hiệu nguyên tử $Z$).
    }
\end{ex}

\begin{ex}
    Các đồng vị của cùng một nguyên tố hóa học là những nguyên tử có
    \choice
    {cùng số neutron nhưng khác nhau về số proton}
    {\True cùng số proton nhưng khác nhau về số neutron, do đó số khối khác nhau}
    {cùng số khối $A$ nhưng khác nhau về số proton}
    {cùng số hạt proton, neutron và electron}
    \loigiai{
        Đồng vị là các nguyên tử của cùng một nguyên tố hóa học, có cùng số proton trong hạt nhân (cùng số hiệu nguyên tử $Z$) nhưng có số neutron ($N$) khác nhau, dẫn tới số khối ($A = Z + N$) của chúng khác nhau.
    }
\end{ex}

\begin{ex}
    Orbital nguyên tử loại $p$ (AO $p$) có dạng hình học nào sau đây?
    \choice
    {Hình cầu}
    {\True Hình số 8 nổi (hình quả tạ đôi)}
    {Hình cánh hoa bốn cánh}
    {Hình elip phẳng}
    \loigiai{
        - Orbital nguyên tử loại $s$ có dạng hình cầu đối xứng quanh hạt nhân.\\
        - Orbital nguyên tử loại $p$ có dạng hình số 8 nổi (hình quả tạ đôi) với ba orbital $p_x, p_y, p_z$ định hướng theo ba trục tọa độ không gian $x, y, z$.
    }
\end{ex}

\begin{ex}
    Các electron thuộc cùng một lớp electron trong nguyên tử có đặc điểm nào sau đây?
    \choice
    {Có mức năng lượng hoàn toàn bằng nhau}
    {\True Có mức năng lượng gần bằng nhau}
    {Có cùng hình dạng orbital nguyên tử}
    {Có cùng chiều tự quay quanh trục}
    \loigiai{
        - Lớp electron gồm các electron có mức năng lượng gần bằng nhau.\\
        - Phân lớp electron gồm các electron có mức năng lượng bằng nhau.
    }
\end{ex}

\begin{ex}
    Trong bảng tuần hoàn các nguyên tố hóa học hiện đại, số thứ tự của chu kì bằng
    \choice
    {số electron ở lớp ngoài cùng của nguyên tử}
    {số electron hóa trị của nguyên tử}
    {\True số lớp electron trong nguyên tử}
    {số hiệu nguyên tử của nguyên tố}
    \loigiai{
        Chu kì là dãy các nguyên tố mà nguyên tử của chúng có cùng số lớp electron, được sắp xếp theo chiều điện tích hạt nhân tăng dần. Do đó, số thứ tự của chu kì luôn bằng số lớp electron trong nguyên tử của nguyên tố đó.
    }
\end{ex}

\begin{ex}
    Trong một chu kì, theo chiều tăng dần của điện tích hạt nhân (từ trái sang phải), tính chất acid -- base của oxide cao nhất và hydroxide tương ứng biến đổi theo quy luật nào sau đây?
    \choice
    {Tính base tăng dần, tính acid giảm dần}
    {\True Tính base giảm dần, tính acid tăng dần}
    {Cả tính acid và tính base đều tăng dần}
    {Cả tính acid và tính base đều giảm dần}
    \loigiai{
        Trong một chu kì, đi từ trái sang phải theo chiều tăng của điện tích hạt nhân, tính kim loại giảm dần và tính phi kim tăng dần, dẫn đến tính base của oxide cao nhất và hydroxide tương ứng giảm dần, đồng thời tính acid của chúng tăng dần.
    }
\end{ex}

\begin{ex}
    Nguyên tố phosphorus ($P$) thuộc nhóm VA trong bảng tuần hoàn. Công thức oxide cao nhất của phosphorus là
    \choice
    {$\text{P}_2\text{O}_3$}
    {$\text{PO}_2$}
    {\True $\text{P}_2\text{O}_5$}
    {$\text{PO}_3$}
    \loigiai{
        Nguyên tố phosphorus ($P$) thuộc nhóm VA, hóa trị cao nhất của nguyên tố nhóm A đối với oxygen bằng số thứ tự của nhóm, tức là bằng $V$. Do đó, công thức oxide cao nhất của phosphorus là $\text{P}_2\text{O}_5$.
    }
\end{ex}

\noindent\textbf{\large PHẦN II. Câu trắc nghiệm đúng sai.} Thí sinh trả lời từ câu 1 đến câu 4. Trong mỗi ý \textbf{a)}, \textbf{b)}, \textbf{c)}, \textbf{d)} ở mỗi câu, thí sinh chọn đúng hoặc sai.

\begin{ex}
    Cho hai nguyên tố $X$ ($Z = 12$) và $Y$ ($Z = 17$).
    \choiceTFt
    {\True Lớp electron thứ 3 ($n = 3$, lớp M) có tối đa 3 phân lớp ($3s, 3p, 3d$) và chứa tối đa 18 electron}
    {\True Cấu hình electron nguyên tử của $X$ là $1s^2 2s^2 2p^6 3s^2$ và của $Y$ là $1s^2 2s^2 2p^6 3s^2 3p^5$}
    {\True Nguyên tử $X$ có 2 electron ở lớp ngoài cùng nên là kim loại; nguyên tử $Y$ có 7 electron ở lớp ngoài cùng nên là phi kim}
    {Để đạt cấu hình electron bền vững của khí hiếm gần nhất, nguyên tử $X$ có xu hướng nhận 6 electron để tạo anion $X^{6-}$}
    \loigiai{
        \begin{itemchoice}
            \itemch Đúng. Lớp $n = 3$ gồm 3 phân lớp: $3s$ (tối đa 2e), $3p$ (tối đa 6e), $3d$ (tối đa 10e). Tổng số electron tối đa của lớp thứ 3 là $2 \times 3^2 = 18$ electron.
            \itemch Đúng. Cấu hình electron nguyên tử:
            - $X$ ($Z = 12$): $1s^2 2s^2 2p^6 3s^2$ (hoặc $[\text{Ne}]3s^2$).
            - $Y$ ($Z = 17$): $1s^2 2s^2 2p^6 3s^2 3p^5$ (hoặc $[\text{Ne}]3s^2 3p^5$).
            \itemch Đúng. Nguyên tử $X$ có 2 electron lớp ngoài cùng ($3s^2$) nên thể hiện tính kim loại điển hình. Nguyên tử $Y$ có 7 electron lớp ngoài cùng ($3s^2 3p^5$) nên thể hiện tính phi kim điển hình.
            \itemch Sai. Nguyên tử $X$ có 2 electron ở lớp ngoài cùng nên có xu hướng nhường đi 2 electron này để đạt cấu hình electron bền vững của khí hiếm neon ($[\text{Ne}]$), tạo thành cation $X^{2+}$, chứ không nhận 6 electron tạo anion $X^{6-}$.
        \end{itemchoice}
    }
\end{ex}

\begin{ex}
    Cho bốn nguyên tố: $_{11}\text{Na}, _{12}\text{Mg}, _{13}\text{Al}$ và $_{19}\text{K}$.
    \choiceTFt
    {\True Các nguyên tố $\text{Na}, \text{Mg}, \text{Al}$ cùng thuộc chu kì 3 trong bảng tuần hoàn vì nguyên tử của chúng đều có 3 lớp electron}
    {Cả bốn nguyên tố $\text{Na}, \text{Mg}, \text{Al}, \text{K}$ đều là các nguyên tố $s$}
    {\True Thứ tự sắp xếp các nguyên tố theo chiều tăng dần của bán kính nguyên tử là: $\text{Al} < \text{Mg} < \text{Na} < \text{K}$}
    {Dung dịch hydroxide tương ứng của potassium ($\text{KOH}$) có tính base yếu hơn dung dịch hydroxide tương ứng của sodium ($\text{NaOH}$)}
    \loigiai{
        \begin{itemchoice}
            \itemch Đúng. Số thứ tự chu kì bằng số lớp electron trong nguyên tử. Cả $\text{Na}, \text{Mg}, \text{Al}$ đều có 3 lớp electron nên đều thuộc chu kì 3.
            \itemch Sai. Cấu hình electron của $\text{Al}$ là $[\text{Ne}]3s^2 3p^1$, có electron cuối cùng điền vào phân lớp $p$ nên aluminium ($\text{Al}$) là nguyên tố $p$, không phải nguyên tố $s$. ($\text{Na}, \text{Mg}, \text{K}$ là nguyên tố $s$).
            \itemch Đúng.
            - Trong cùng chu kì 3: theo chiều tăng dần của điện tích hạt nhân, bán kính nguyên tử giảm dần: $r_{\text{Al}} < r_{\text{Mg}} < r_{\text{Na}}$.
            - So sánh $\text{Na}$ (chu kì 3, có 3 lớp electron) và $\text{K}$ (chu kì 4, có 4 lớp electron): $\text{K}$ có nhiều lớp electron hơn nên $r_{\text{Na}} < r_{\text{K}}$.
            - Do đó thứ tự tăng dần của bán kính nguyên tử là: $\text{Al} < \text{Mg} < \text{Na} < \text{K}$.
            \itemch Sai. Potassium ($\text{K}$) và sodium ($\text{Na}$) cùng thuộc nhóm IA. Trong một nhóm A, theo chiều từ trên xuống dưới tính kim loại tăng dần ($\text{K} > \text{Na}$), do đó tính base của hydroxide tương ứng tăng dần: tính base của $\text{KOH}$ mạnh hơn $\text{NaOH}$.
        \end{itemchoice}
    }
\end{ex}

\begin{ex}
    Cho hai nguyên tố phi kim thuộc chu kì 3 là phosphorus ($Z = 15$) và sulfur ($Z = 16$).
    \choiceTFt
    {\True Cấu hình electron nguyên tử sulfur là $[\text{Ne}]3s^2 3p^4$, nguyên tử có 6 electron hóa trị và thuộc nhóm VIA}
    {\True Cả phosphorus và sulfur đều thuộc khối các nguyên tố $p$ trong bảng tuần hoàn}
    {\True Hydroxide tương ứng với oxide cao nhất của sulfur là $\text{H}_2\text{SO}_4$, là một acid mạnh}
    {Trong cùng chu kì 3, tính acid của phosphoric acid ($\text{H}_3\text{PO}_4$) mạnh hơn tính acid của sulfuric acid ($\text{H}_2\text{SO}_4$)}
    \loigiai{
        \begin{itemchoice}
            \itemch Đúng. Sulfur có $Z = 16$, cấu hình electron là $1s^2 2s^2 2p^6 3s^2 3p^4$ (hay $[\text{Ne}]3s^2 3p^4$), có 6 electron hóa trị ở lớp ngoài cùng, thuộc chu kì 3, nhóm VIA.
            \itemch Đúng. Cả hai nguyên tử đều có electron cuối cùng điền vào phân lớp $p$ ($3p^3$ ở P và $3p^4$ ở S) nên đều là các nguyên tố $p$.
            \itemch Đúng. Hóa trị cao nhất của sulfur là VI, oxide cao nhất là $\text{SO}_3$, hydroxide tương ứng là $\text{H}_2\text{SO}_4$, đây là một acid vô cơ mạnh.
            \itemch Sai. Trong chu kì 3, đi từ trái sang phải theo chiều tăng điện tích hạt nhân ($Z = 15 \rightarrow 16$), tính phi kim tăng dần nên tính acid của hydroxide tương ứng tăng dần: $\text{H}_3\text{PO}_4$ có tính acid yếu hơn $\text{H}_2\text{SO}_4$.
        \end{itemchoice}
    }
\end{ex}

\begin{ex}
    Xét quy luật biến đổi tuần hoàn tính chất của các nguyên tố hóa học trong bảng tuần hoàn.
    \choiceTFt
    {\True Tính chất của các nguyên tố và đơn chất biến đổi tuần hoàn theo chiều tăng dần của điện tích hạt nhân nguyên tử}
    {\True Trong một chu kì, đi từ trái sang phải, bán kính nguyên tử giảm dần do điện tích hạt nhân tăng trong khi số lớp electron giữ nguyên}
    {Trong một nhóm A, đi từ trên xuống dưới, độ âm điện của các nguyên tử tăng dần}
    {\True Quy luật biến đổi tuần hoàn tính chất của các nguyên tố thực chất bắt nguồn từ sự lặp lại tuần hoàn cấu hình electron lớp ngoài cùng của nguyên tử}
    \loigiai{
        \begin{itemchoice}
            \itemch Đúng. Đây là nội dung cơ bản của định luật tuần hoàn các nguyên tố hóa học hiện đại.
            \itemch Đúng. Trong cùng chu kì, số lớp electron bằng nhau, nhưng điện tích hạt nhân $Z$ tăng dần làm tăng lực hút tĩnh điện giữa hạt nhân và electron, kéo các electron lại gần hạt nhân hơn, làm bán kính nguyên tử giảm dần.
            \itemch Sai. Trong một nhóm A, từ trên xuống dưới số lớp electron tăng, bán kính nguyên tử tăng nhanh làm giảm lực hút của hạt nhân đối với electron liên kết, dẫn tới giá trị độ âm điện giảm dần.
            \itemch Đúng. Cấu hình electron lớp ngoài cùng quyết định tính chất hóa học. Sự lặp lại tuần hoàn cấu hình electron lớp ngoài cùng sau mỗi chu kì chính là nguyên nhân sâu xa dẫn đến sự biến đổi tuần hoàn tính chất của các đơn chất và hợp chất.
        \end{itemchoice}
    }
\end{ex}

\noindent\textbf{\large PHẦN III. Câu trắc nghiệm trả lời ngắn.} Thí sinh trả lời từ câu 1 đến câu 4.

\begin{ex}
    Cho các nguyên tố có số hiệu nguyên tử lần lượt là $X$ ($Z = 11$), $Y$ ($Z = 12$), $Z$ ($Z = 15$), $T$ ($Z = 16$), $Q$ ($Z = 17$). Dựa vào đặc điểm cấu hình electron lớp ngoài cùng của các nguyên tử trên, hãy cho biết có bao nhiêu nguyên tố thể hiện tính chất hóa học cơ bản là phi kim?
    \shortans{3}
    \loigiai{
        Viết cấu hình electron của các nguyên tử và xác định số electron lớp ngoài cùng:
        \begin{itemize}
            \item $X$ ($Z = 11$): $1s^2 2s^2 2p^6 3s^1 \rightarrow$ có 1 electron lớp ngoài cùng $\rightarrow$ là kim loại.
            \item $Y$ ($Z = 12$): $1s^2 2s^2 2p^6 3s^2 \rightarrow$ có 2 electron lớp ngoài cùng $\rightarrow$ là kim loại.
            \item $Z$ ($Z = 15$): $1s^2 2s^2 2p^6 3s^2 3p^3 \rightarrow$ có 5 electron lớp ngoài cùng $\rightarrow$ là phi kim.
            \item $T$ ($Z = 16$): $1s^2 2s^2 2p^6 3s^2 3p^4 \rightarrow$ có 6 electron lớp ngoài cùng $\rightarrow$ là phi kim.
            \item $Q$ ($Z = 17$): $1s^2 2s^2 2p^6 3s^2 3p^5 \rightarrow$ có 7 electron lớp ngoài cùng $\rightarrow$ là phi kim.
        \end{itemize}
        Vậy trong 5 nguyên tố đã cho, có 3 nguyên tố thể hiện tính chất phi kim là $Z, T, Q$.
    }
\end{ex}

\begin{ex}
    Một nguyên tử của nguyên tố $R$ có tổng số các hạt cơ bản (proton, neutron, electron) là $52$. Trong nguyên tử $R$, tổng số hạt mang điện nhiều hơn số hạt không mang điện là $16$. Dựa vào cấu hình electron nguyên tử, hãy xác định số thứ tự chu kì của nguyên tố $R$ trong bảng tuần hoàn các nguyên tố hóa học.
    \shortans{3}
    \loigiai{
        Gọi $p, n, e$ lần lượt là số hạt proton, neutron, electron trong nguyên tử $R$. Do nguyên tử trung hòa về điện nên $p = e$.\\
        Theo đề bài, ta có hệ phương trình:
        \[
            \begin{cases}
                2p + n = 52\\
                2p - n = 16
            \end{cases}
            \iff
            \begin{cases}
                4p = 68\\
                n = 2p - 16
            \end{cases}
            \iff
            \begin{cases}
                p = 17\\
                n = 18.
            \end{cases}
        \]
        Số hiệu nguyên tử của $R$ là $Z = p = 17$.\\
        Cấu hình electron của nguyên tử $R$ là: $1s^2 2s^2 2p^6 3s^2 3p^5$.\\
        Vì nguyên tử $R$ có 3 lớp electron ($n = 1, 2, 3$) nên nguyên tố $R$ thuộc chu kì 3 trong bảng tuần hoàn.
    }
\end{ex}

\begin{ex}
    Cho các nguyên tố thuộc chu kì 3 gồm: $_{11}\text{Na}, _{12}\text{Mg}, _{13}\text{Al}, _{14}\text{Si}, _{15}\text{P}, _{16}\text{S}, _{17}\text{Cl}$. Dựa vào xu hướng biến đổi tuần hoàn bán kính nguyên tử trong một chu kì, hãy cho biết có bao nhiêu nguyên tố có bán kính nguyên tử nhỏ hơn bán kính nguyên tử của phosphorus ($_{15}\text{P}$)?
    \shortans{2}
    \loigiai{
        Trong một chu kì, đi từ trái sang phải theo chiều tăng dần của điện tích hạt nhân:
        - Số lớp electron của các nguyên tử không đổi (đều có 3 lớp electron).
        - Điện tích hạt nhân $Z$ tăng dần làm cho lực hút tĩnh điện giữa hạt nhân với các electron lớp ngoài cùng tăng dần, kéo các electron lại gần hạt nhân hơn.
        - Do đó, bán kính nguyên tử giảm dần từ trái sang phải:
        \[
            r_{\text{Na}} > r_{\text{Mg}} > r_{\text{Al}} > r_{\text{Si}} > r_{\text{P}} > r_{\text{S}} > r_{\text{Cl}}.
        \]
        Các nguyên tố có bán kính nguyên tử nhỏ hơn bán kính nguyên tử của phosphorus ($_{15}\text{P}$) là sulfur ($_{16}\text{S}$) và chlorine ($_{17}\text{Cl}$).\\
        Vậy có tất cả 2 nguyên tố thỏa mãn.
    }
\end{ex}

\begin{ex}
    Nguyên tố $X$ thuộc chu kì 3, nhóm VA trong bảng tuần hoàn các nguyên tố hóa học. Cho các nhận định sau về tính chất của $X$ và các hợp chất tạo nên từ $X$:
    \begin{itemize}
        \item[(1)] Đơn chất $X$ thể hiện tính chất hóa học cơ bản là tính phi kim.
        \item[(2)] Hóa trị cao nhất của nguyên tố $X$ đối với oxygen bằng $5$.
        \item[(3)] Công thức oxide cao nhất của $X$ là $\text{X}_2\text{O}_5$ và có tính acid (acidic oxide).
        \item[(4)] Hydroxide ứng với hóa trị cao nhất của $X$ có công thức $\text{H}_3\text{XO}_4$ và có tính base mạnh.
    \end{itemize}
    Có bao nhiêu nhận định \textbf{đúng} trong số các nhận định trên?
    \shortans{3}
    \loigiai{
        Nguyên tố $X$ thuộc chu kì 3, nhóm VA ($X$ là phosphorus, $Z = 15$):
        \begin{itemize}
            \item[(1)] Đúng. Cấu hình electron nguyên tử là $1s^2 2s^2 2p^6 3s^2 3p^3$, có 5 electron lớp ngoài cùng nên đơn chất $X$ thể hiện tính chất phi kim.
            \item[(2)] Đúng. Nguyên tố thuộc nhóm VA nên có hóa trị cao nhất với oxygen bằng số thứ tự của nhóm, tức là bằng 5.
            \item[(3)] Đúng. Công thức oxide cao nhất là $\text{X}_2\text{O}_5$ ($\text{P}_2\text{O}_5$), là một acidic oxide, tác dụng với nước tạo thành acid tương ứng.
            \item[(4)] Sai. Hydroxide cao nhất tương ứng của $X$ là $\text{H}_3\text{XO}_4$ ($\text{H}_3\text{PO}_4$, phosphoric acid) là một acid có độ mạnh trung bình, không phải là chất có tính base mạnh.
        \end{itemize}
        Vậy có 3 nhận định đúng là (1), (2) và (3).
    }
\end{ex}

\noindent\textbf{\large PHẦN IV. Tự luận.} Thí sinh trả lời từ câu 1 đến câu 2.

\begin{ex}
    Nguyên tử của một nguyên tố phi kim $M$ có tổng số electron ở các phân lớp $p$ bằng $11$.
    \begin{itemize}
        \item[a)] Viết cấu hình electron đầy đủ của nguyên tử $M$.
        \item[b)] Xác định vị trí của nguyên tố $M$ trong bảng tuần hoàn (ô nguyên tố, chu kì, nhóm) và giải thích ngắn gọn.
        \item[c)] Viết công thức oxide cao nhất và hydroxide tương ứng của $M$. Nêu tính chất acid -- base của các hợp chất này.
    \end{itemize}
    \loigiai{
        \begin{itemize}
            \item[a)] Tổng số electron ở các phân lớp $p$ là 11, suy ra các phân lớp $p$ lần lượt là $2p^6$ và $3p^5$.\\
            Cấu hình electron đầy đủ của nguyên tử $M$ là:
            \[
                1s^2 2s^2 2p^6 3s^2 3p^5 \quad (\text{hoặc } [\text{Ne}]3s^2 3p^5).
            \]
            \item[b)] Vị trí của $M$ trong bảng tuần hoàn:
            \begin{itemize}
                \item Ô nguyên tố: Số 17 (do số hiệu nguyên tử $Z = 17$, nguyên tử có 17 electron).
                \item Chu kì: Chu kì 3 (do nguyên tử $M$ có 3 lớp electron).
                \item Nhóm: Nhóm VIIA (do có electron cuối cùng điền vào phân lớp $p$ và có 7 electron hóa trị ở lớp ngoài cùng $3s^2 3p^5$).
            \end{itemize}
            (Nguyên tố $M$ là chlorine, kí hiệu Cl).
            \item[c)] Công thức oxide cao nhất và hydroxide tương ứng:
            \begin{itemize}
                \item Oxide cao nhất: $\text{Cl}_2\text{O}_7$ (dichlorine heptoxide), là một acidic oxide (oxit axit mạnh).
                \item Hydroxide tương ứng: $\text{HClO}_4$ (perchloric acid), là một acid rất mạnh.
            \end{itemize}
        \end{itemize}
    }
\end{ex}

\begin{ex}
    Hai nguyên tố $A$ và $B$ đứng kế tiếp nhau trong cùng một chu kì của bảng tuần hoàn (thuộc các nhóm A). Tổng số hạt mang điện trong cả hai hạt nhân nguyên tử của $A$ và $B$ là $25$ ($Z_A < Z_B$).
    \begin{itemize}
        \item[a)] Xác định số hiệu nguyên tử và tên gọi của hai nguyên tố $A$ và $B$.
        \item[b)] Viết cấu hình electron nguyên tử của $A$ và $B$.
        \item[c)] So sánh bán kính nguyên tử và độ âm điện giữa $A$ và $B$. Giải thích ngắn gọn dựa trên cấu tạo nguyên tử.
    \end{itemize}
    \loigiai{
        \begin{itemize}
            \item[a)] Hạt mang điện trong hạt nhân là proton. Tổng số hạt proton của hai nguyên tử là $25$ và hai nguyên tố kế tiếp nhau trong cùng chu kì ($Z_B - Z_A = 1$):
            \[
                \begin{cases}
                    Z_A + Z_B = 25\\
                    Z_B - Z_A = 1
                \end{cases}
                \iff
                \begin{cases}
                    Z_A = 12\\
                    Z_B = 13.
                \end{cases}
            \]
            Vậy:
            - Nguyên tố $A$ có $Z_A = 12$ là magnesium (kí hiệu Mg).
            - Nguyên tố $B$ có $Z_B = 13$ là aluminium (kí hiệu Al).
            \item[b)] Cấu hình electron nguyên tử:
            - $A$ ($Z = 12$): $1s^2 2s^2 2p^6 3s^2$ (hoặc $[\text{Ne}]3s^2$).
            - $B$ ($Z = 13$): $1s^2 2s^2 2p^6 3s^2 3p^1$ (hoặc $[\text{Ne}]3s^2 3p^1$).
            \item[c)] So sánh bán kính nguyên tử và độ âm điện:
            - Cả $A$ (Mg) và $B$ (Al) đều thuộc chu kì 3, nguyên tử đều có 3 lớp electron.
            - Do điện tích hạt nhân của $B$ lớn hơn $A$ ($13+ > 12+$) nên lực hút tĩnh điện giữa hạt nhân với các electron lớp ngoài cùng của $B$ mạnh hơn, kéo electron lại gần hạt nhân hơn.\\
            $\rightarrow$ Bán kính nguyên tử: $r_{\text{Mg}} > r_{\text{Al}}$ (hay $r_A > r_B$).\\
            $\rightarrow$ Độ âm điện: $\chi_{\text{Mg}} < \chi_{\text{Al}}$ (hay $\chi_A < \chi_B$).
        \end{itemize}
    }
\end{ex}
